\section{Limitations} \label{sec:limitations}

Despite our attempt to devise novel counterfactual conditions to gauge an LM's ``true'' reasoning ability, it may not be precisely reflected by the counterfactual performance due to several factors.

\subsection{Underestimation} \label{sec:underestimation}

For our main evaluations, we aim to construct counterfactual tasks that have the same difficulty as the default variants so that task difficulty does not confound our comparisons. This is not always possible---in fact, an objective difficulty measure may not even exist. One could, for example, argue that base-11 addition is harder than base-10 because it requires reasoning with one additional digit, or base-9 is harder than base-10 because on average the sums would consist of more digits. 



Retrieving notes in melodies in different keys faces a similar issue.
We expect similar retrieval difficulty under different keys
if the model recalls a melody as a series of abstract relations in a scale and directly maps them onto notes in a target key.
However, an alternative strategy would be to first retrieve the note in a canonical key and then \emph{transpose} it to the desired uncommon key. This 2-step process is a natural one that is often employed by musicians. And with this strategy, the counterfactual task consists of 2 steps and is harder than (and requires first) completing the 1-step original task. The counterfactual setup thus introduces a confounder: low performance may be driven by the increased difficulty of the counterfactual task, rather than overfitting to melodies in their canonical keys, \emph{if} models are employing two-step strategy. However, since both strategies are available to models and we do not prompt them to use a particular one, reliance on this two-step strategy may itself be indicative of overfitting to the original canonical keys.

\subsection{Overestimation} \label{sec:overestimation}

We can never be certain of how rarely particular counterfactual conditions are encountered during pretraining. It is quite likely that there is text online that, for example, draws rotated versions of various objects used in our study. Consequently, the effect of overfitting could also manifest in our counterfactual conditions, and the default-counterfactual gap could actually be larger for some genuinely unseen conditions.

We also distinguish between two types of counterfactual perturbations. One type fundamentally affects the operation of the world model and necessitates an understanding of the counterfactual world to perform the task in it (e.g., arithmetic base or 1-based indexing\footnote{It may be tempting to consider a simple replacement strategy \texttt{[i]}\textrightarrow\texttt{[i-1]} to map back to 0-based indexing. But this does not work for the dictionary type. There are other complications; see Table~\ref{tab:prompts-programming}.}). On the other hand, some perturbations are more superficial and may admit a shortcut where the model first figures out a simple mapping of the input back to the default conditions and performs the task (potentially leveraging instance-level memorization) under those.
In some of our tasks, this mapping may be simple, such as the word replacements in the natural language logical reasoning task\footnote{To be more concrete, imagine that a model memorizes an instance with nine premises on dogs involving complex logical relationships, and that it entails a given conclusion. For the counterfactual instance, we replace the word ``dogs'' with another object, say ``headphones,'' to make the premises no longer factually true. Instead of performing the reasoning over premises with headphones such as how they are, counterfactually, the cutest creatures, a model could identify the mapping ``dogs'' $\to$ ``headphones'', revert it (i.e., replace all ``headphones'' back to ``dogs''), and perform the task under the default common-sense-complying conditions.} (\S\ref{sec:folio-task}) and the transformation functions for the drawing task (\S\ref{sec:drawing-task}), which could potentially be exploited by the models. We explicitly disallow this in our prompt for the drawing task (Table~\ref{tab:prompts-drawing}) but did not identify a good way to forbid this for logical reasoning, potentially accounting for its generally high counterfactual performance.

Finally, we reiterate from \S\ref{sec:results} that a non-perfect CCC accuracy does not allow us to perfectly tease apart counterfactual performance and a failure of counterfactual condition comprehension. But often the default-counterfactual gap is so prominent that it is still strongly suggestive of overfitting to the default conditions. Also, recall from \S\ref{sec:counterfactual-eval} that the CCC itself is also a nontrivial task. For ThonPy, for example, the CCC also involves program evaluation, albeit with simpler statements that involve less reasoning, such as \texttt{print("qrstu"[4])}. We do not see an easy way to introduce ThonPy CCC that is entirely disentangled from program evaluation. This conflation would result in the CCC accuracy's being lower than what would reflect the model's understanding of the counterfactual conditions.
